\renewcommand{\bibname}{Quellen und Bandzuordnung}
\begin{thebibliography}{9}
\addcontentsline{toc}{chapter}{Quellen und Bandzuordnung}
\footnotesize
\setlength{\itemsep}{4pt}
\setlength{\parskip}{0pt}
\interlinepenalty=10000

\bibitem{project}
Martin Kunze: \emph{Die Grundlagen der Mathematik}.
Band~4, Kapitel~2: Totalitätsaxiom, Definition der totalen Relation
und Mitgliedschaftsrelation.
Band~3: geordnete Paare, kartesische Produkte, Aussonderung,
Potenzmengen und Mengenoperationen.
Die Mitgliedschaftsrelation wird hier unter der genauen Bedingung
$\varnothing\notin\mathcal M$ aufgeführt; die in Band~4 verwendete
Disjunktheitsannahme ist dafür entbehrlich.
Die strikte Obermengenrelation aus Band~4 wird hier nicht vertieft,
um die Ordnungstheorie ihren eigenen Bänden zu belassen.
Format nach den Leitfäden zu Band~3, Band~28 und Band~37.

\bibitem{meseguer}
José Meseguer: \emph{Set Theory and Algebra in Computer Science:
A Gentle Introduction to Mathematical Modeling}.
University of Illinois at Urbana-Champaign, Fassung vom
29.~August 2024.
Abschnitt~5.3, Druckseiten~32--33: erreichte Mengen und
Totalität; Fußnote~5: Abgrenzung von Vergleichbarkeit.
Abschnitt~5.4, Druckseite~33: mengentheoretische Definition
der Verkettung.
\url{https://courses.physics.illinois.edu/cs476/fa2024/readings/STACS.pdf}.

\end{thebibliography}

\noindent\textbf{Einordnung der Zusammenstellung.}
Die Kriterien zu Einschränkungen, Vereinigungen, Durchschnitten,
Differenzen und Paarbildungen sind elementare Folgerungen aus
der Totalitätsdefinition und den Mengenoperationen.
Sie sind hier als zusammenhängende Referenzliste formuliert;
es werden weder neue Forschungsergebnisse noch historische
Erstentdeckungen beansprucht.
Die externe Quelle dient der Begriffs- und Notationskontrolle,
nicht als Behauptung, sämtliche hier formulierten Sätze stünden
dort in dieser Gestalt.
